\def\ChapterCode{AUP}
\chapter
[Anamnese und Patientengespräch]
{Anamnese und Patientengespräch}
\label{chp:AUU}
\label{chp:AUP}

\ChapterIntro
{%
\noindent Dieses Kapitel behandelt die Durchführung eines
\E{Patientengespräches} und der Erhebung der \E{Anamnese}.
Beides sind wichtige Basismaßnahmen um eine (Verdachts-)
Diagnose erarbeiten zu können. Das Patientengespräch ist
außerdem ein wichtiger Bestandteil der \E{psychischen
Betreuung}.


}
{%
\begin{description}
  \item[Maintainer:] Sebastian Gabriel
  \item[Autoren:] Diverse
  \item[Reviewer:] --
  \item[Version:] {\DocVersion} ({\DocVersionInfo})
  \item[SHA1:] {\DocGitCommit}
\end{description}

{\TxTeilkapitelAASS}
}

% \begin{WidePar}
% \Parallel{
% \minisec{{}}
% 
% }{
% 
% }
% \end{WidePar}

\protect{\begin{verbatim}
 ________________________________________
/ You mentioned your name as if I should \
| recognize it, but beyond the obvious   |
| facts that you are a bachelor, a       |
| solicitor, a freemason, and an         |
| asthmatic, I know nothing whatever     |
| about you.                             |
|                                        |
| -- Sherlock Holmes, "The Norwood       |
\ Builder"                               /
 ----------------------------------------
       \   ,__,
        \  (oo)____
           (__)    )\
              ||--|| *

\end{verbatim}}

\clearpage
% \TextSlide{\TxQuerverweise}
% {~
% 
% \noindent\includegraphics[angle=270,width=\linewidth]{./images/leitsystem/Leitschema-PAM_AUP.pdf}
% }{\begin{itemize}
%   \item \EE{Medizinprodukte allgemein}: \CO{[MP1]} (\ref{chp:MP1})
%   \item \EE{Untersuchungen}:                 \CO{[BAU]} (\ref{chp:BAU})
%   \item \EE{Sanitätshilfemaßnahmen}:         \CO{[SAN]} (\ref{chp:SAN})
% %   \item \EE{Anamnese und Patientengespräch}: \CO{[AUP]} (\ref{chp:AUP})
% \end{itemize}
% }{}{}{}{}




\Layout{2}{Das Patientengespräch}{
Als Patientengespräch versteht man das professionelle
Gespräch mit dem Patienten sowie beteiligten Dritten
(Angehörige, Beobachter, \ldots ) zum Zweck der Klärung der
Situation, Erhebung der Anamnese und der psychischen
Betreuung.
}{}{}{}{}


%========Beginn Abschnitt Emhofer============================================



% \section{Die Anamnese erz\"ahlt uns die
% Geschichte des Patienten}
\index{Anamnese}

\section{Kommunikationsregeln}
\label{topic:kommunikationsregeln}

% \begin{wrapfigure}{r}{4cm}
% \includegraphics[width=\linewidth]{./images/teaser/imgp0482_00800.jpg}
% \end{wrapfigure}
Der Erfolg des Patienten- bzw. Anamnesegesprächs ist für das
weitere Vorgehen entscheidend. Oft erhält man wichtige
Information über den Zustand des Patienten nicht durch die
gemessenen Werte, sondern durch eine professionelle und
zielgerichtete Befragung. Neben dem Informationsgewinn wirkt dies zusätzlich
vertrauensfördernd.
% Der Erfolg des Gespräches hängt dabei einerseits vom Geschick des Gesprächführenden,
% andererseits aber auch von der Bereitschaft des Patienten
% (oder dessen Umgebung!) ab.

\TextSlide{Grundregeln und Tipps}{
Ein gut geführtes Patientengespräch ist durch folgende Merkmale gekennzeichnet:
\begin{itemize}
    \item \EE{Vorstellen}: Zur Vorstellung gehört die Nennung des eigenen
    Namens und -- so es die Dringlichkeit zulässt -- ein Händedruck.
    \item \EE{Augenhöhe}: Geht man in die Hocke kann man mit dem Patient auf
    \E{Augenhöhe} kommunizieren. Dies nimmt dem Patient Angst und fördert eine
    Gesprächsbasis \E{auf einer Ebene}.
    \item \EE{Deutlichkeit und Einfachheit}: Eine deutliche und einfache
    Ausdrucksweise ist besonders wichtig um Missverständnisse zu vermeiden und
    schnell zur gewünschten Information zu kommen.\footnote{Evtl.
    benötigt der Patient Hörhilfen.
    Diese müssen evtl. erst eingeschalten werden.} Fachbegriffe sind zu
    vermeiden!
    \item \EE{Aktives Zuhören}: Die
    Antworten des Patienten müssen beim Helfer \enquote{ankommen} und nicht
    \enquote{verloren gehen}. Wichtige Informationen werden oft
    nur einmal berichtet, der Patient nimmt im weiteren
    Verlauf an, dass er sie bereits gesagt und das
    Fachpersonal sie verstanden hat.
    
    Negativbeispiel: \begin{quote}\begin{inparadesc}
                     \item[Sanitäter:] \Speech{\enquote{Haben Sie Schmerzen?}}
                     \item[Patient:] \Speech{\enquote{Nein.}}
                     \item[Sanitäter:] \Speech{\enquote{Wie stark ist der Schmerz?}}
                     \end{inparadesc}\end{quote}
    \item \EE{Verstehen}: Es ist wichtig zu verstehen was der
    Patient mit dem Gesagten \E{meint}! Um sich zu vergewissern, dass man die
    Information richtig verarbeitet hat, kann man das Gesagte in
    eigenen Worten wiederholen und sich vom Patient die Richtigkeit bestätigen
    lassen.
    \item \EE{\enquote{Monogam}}: Es ist immer nur \E{eine
    Frage gleichzeitig} zu stellen.
    \item \EE{Information}: Der Patient ist immer über das Vorgehen der Helfer
    zu informieren!
    \item \EE{Aufrichtigkeit}: Die Fragen des
    Patienten sind nach bestem Wissen zu beantworten! \EE{Das Anlügen des
    Patienten ist \E{tabu}!} Liegen mögliche Antworten außerhalb des Wissens-
    oder Kompetenzbereiches ist auf ärztlichen Rat zu verweisen.
    \item \EE{Bezugsperson}: Sind mehrere Helfer
    vorort, so soll nur einer der Ansprechpartner des
    Patienten sein!\label{topic:einer-redet}
\end{itemize}
}{
\begin{itemize}
  \item Vorstellen
  \item auf Augenhöhe mit dem Patienten
  \item deutliche, einfache Sprache
  \item aktives Zuhören
  \item Verstehen der Antworten
  \item nur eine Frage gleichzeitig
  \item Informieren über die Maßnahmen
  \item Aufrichtigkeit
  \item Bezugsperson
\end{itemize}
}{}{}{}

% \TextSlide{Grundregeln und Tipps}{
% Im Folgenden möchten wir Ihnen einige \E{Grundregeln} und
% wertvolle Tipps geben, wie Sie Ihr Patientengespräch so
% gestalten können, dass Sie die gewünschten Informationen
% erhalten und ein Vertrauensverhältnis mit dem Patienten
% aufbauen können:
% 
% \begin{itemize}
%     \item \EE{Vorstellen}: Stellen Sie sich mit Ihrem Namen
%     und Ihrer Funktion vor und - so es der Notfall zulässt -
%     geben Sie dem Patienten die Hand!
%     \item \EE{Augenhöhe}: Begeben Sie sich auf die gleiche
%     Ebene (auf \E{Augenhöhe}) mit dem Patienten!
%     \item \EE{Deutlich} Sprechen Sie laut und
%     deutlich!\footnote{Evtl. benötigt der Patient Hörhilfen.
%     Diese müssen evtl. erst eingeschalten werden.}
%     \item \EE{Aktives Zuhören}: Die
%     Antworten sollen bei Ihnen \enquote{ankommen} und nicht
%     \enquote{verloren gehen}. Wichtige Informationen werden oft
%     nur einmal berichtet, der Patient nimmt im weiteren
%     Verlauf an, dass er sie bereits gesagt und das
%     Fachpersonal sie verstanden hat.
%     
%     Negativbeispiel: \begin{quote}\begin{inparadesc}
%                      \item[Sanitäter:] \Speech{\enquote{Haben Sie Schmerzen?}}
%                      \item[Patient:] \Speech{\enquote{Nein.}}
%                      \item[Sanitäter:] \Speech{\enquote{Wie stark ist der Schmerz?}}
%                      \end{inparadesc}\end{quote}
%     \item \EE{Verstehen}: Versuchen Sie zu verstehen was der
%     Patient \E{meint}! Wiederholen Sie wichtige
%     Informationen um sicher zu gehen dass Sie es richtig
%     verstanden haben.
%     \item \EE{Einfachheit}: Verwenden Sie eine leicht
%     verständliche Sprache und vermeiden Sie Fachausdrücke!
%     \item \EE{\enquote{Monogam}}: Stellen Sie immer nur \E{eine
%     Frage gleichzeitig}.
%     \item \EE{Information}: Informieren Sie den Patienten über
%     Ihr Vorgehen!
%     \item \EE{Aufrichtigkeit}: Versuchen Sie, Fragen des
%     Patienten nach bestem Wissen zu beantworten, aber
%     erfinden Sie nichts! \EE{Lügen Sie den Patienten
%     \E{nicht} an!} Treffen Sie keine Aussagen oder
%     Vorhersagen, die außerhalb Ihres Kompetenz- und
%     Zuständigkeitsbereichs liegen.
%     \item \EE{Bezugsperson}: Im Team sollten nicht alle auf
%     den Patienten einreden! Der Patient sollte wenn möglich nur
%     einen Ansprechpartner
%     haben!\label{topic:einer-redet}
% \end{itemize}
% }{
% \begin{itemize}
%   \item Vorstellen
%   \item auf Augenhöhe mit dem Patienten
%   \item deutliche, einfache Sprache
%   \item aktives Zuhören
%   \item Verstehen der Antworten
%   \item nur eine Frage gleichzeitig
%   \item Informieren über die Maßnahmen
%   \item Aufrichtigkeit
%   \item Bezugsperson
% \end{itemize}
% }{}{}{}

\TextSlide{Fragentypen}
{
Grundsätzlich unterscheidet man \EE{offene} und
\EE{geschlossene}, sowie \EE{gezielte} Fragen und
\EE{Suggestivfragen}.

\begin{itemize}
  \item \EE{Offene Fragen} kann der Patient frei, mit
  seinen eignen Worten beantworten. Je nach Patient kann man so
  mit wenig Aufwand sehr viele Informationen gewinnen. Das
  Problem dabei kann jedoch sein, dass die Informationen
  ungeordnet und mitunter auch in der Situation mäßig wichtig
  sind.
  
  \begin{quote}
    \Speech{\enquote{Was haben Sie für ein Problem?}}
    
    \Speech{\enquote{Bitte beschreiben Sie den Schmerz.}}
  \end{quote}
  
  
  \item \EE{Geschlossene Fragen} sind sehr direkt und
  speziell, der Patient kann i.\,d.\,R. nur mit
  \emph{\enquote{ja}}, \emph{\enquote{nein}} oder
  \emph{\enquote{weiß nicht}} antworten. Sie machen daher
  eher nur dann Sinn, wenn man ganz spezielle Dinge bewusst
  direkt abfragen kann oder muss. \Ca{Achtung!} Manche
  Patienten neigen dazu, grundsätzlich \enquote{ja} zu
  sagen, ohne dass sie die Frage wirklich verstanden haben!
  
  \begin{quote}
    \Speech{\enquote{Haben Sie sich vorher angestrengt?}}
    
    \Speech{\enquote{Haben Sie das schon einmal gehabt?}}
  \end{quote}
  
  \item \EE{Gezielte Fragen} fokussieren die Antwort des
  Patienten und
  lassen sich nur schwer \enquote{umgehen}. Sie dienen der
  Präzisierung. Werden zu
  viele von diesen Fragen gestellt, kann das Gespräch leicht
  in eine \enquote{Einbahnstraße} abgleiten. Gezielte Fragen können
  sowohl offen, als auch geschlossen sein.
  \begin{quote}
    \Speech{\enquote{Wohin strahlt der
    Schmerz aus?}}
  \end{quote}
  
  \item \EE{Suggestivfragen} sind \enquote{böse}: Hier
  legt der Gesprächsführende dem Patienten die Antworten \enquote{in
  den Mund}. Sie können sowohl als offene, als auch als
  geschlossene Fragen formuliert werden.
  
  \begin{quote}
    \Speech{\enquote{Haben Sie \emph{eh keine} Schmerzen?}}
    
    \Speech{\enquote{Ist der Schmerz stechend?}}
  \end{quote}
  
  Dieser Fragentyp soll nicht -- oder nur im begründeten
  Ausnahmefall -- angewendet werden.
\end{itemize}



}
{
\begin{itemize}
    \item Offene Fragen
    \item Geschlossene Fragen
    \item Gezielte Fragen
    \item Suggestivfragen 
\end{itemize}
}{}{}{}

\TextSlide{Kinder}
{Die Kommunikation mit Kindern kann besonders herausfordernd
sein. Der Kontakt mit Einrichtungen des Gesundheitswesens,
sei es in einer Ordination, im Spital oder im Rahmen des
Rettungsdienstes, stellt eine aufregende und i.\,d.\,R.
angsteinflößende Ausnahmesituation dar. Die folgenden
Vorschläge können unter Umständen hilfreich sein -- ein
allgemeingültiges \enquote{Kochrezept} gibt es allerdings nicht:

\begin{itemize}
  \item \EE{Kleine Erwachsene}: Manchmal sind Kinder doch
  einfach \enquote{kleine Erwachsene}: Die elementaren Grundsätze
  \E{Aufrichtigkeit} und \E{Aufklärung}, welche für das
  Verhalten gegenüber Erwachsenen gilt, gilt genauso für
  Kinder. Ein altersangepasstes Auftreten ist zwar oft
  notwendig, darf aber nicht als Ausrede benutzt werden, um
  die Grundsätze zu vernachlässigen.
%   \item \EE{Aufrichtigkeit}: Seien Sie immer ehrlich! Kinder
%   merken sehr schnell, wenn Sie lügen. Lügen  Sie ein Kind
%   an, so wird das Kind schnell das Vertrauen in Sie
%   verlieren  und jede weitere Untersuchung verweigern!
% 
%   Kinder nehmen non-verbale Signale stärker wahr als
%   Erwachsene. Versuchen Sie daher durch einen freundlichen
%   Gesichtsausdruck und eine offene Körperhaltung Vertrauen
%   zu bilden!
  \item \EE{Aufrichtigkeit}: Kinder bemerken sehr schnell, wenn man lügt. Ein
  ehrlicher Umgang mit Kindern ist daher besonders wichtig, da sonst das
  Vertrauen verloren ist und das Kind möglicherweise spätere Untersuchungen
  (aus Angst) verweigert.
%   \item \EE{Erklären} Sie dem Kind (altersgerecht) warum Sie
%   mit ihm oder ihr reden oder es untersuchen möchten.
  \item \EE{Erklären}: Das (altersgerechte) Erklären von Untersuchungen und
  Maßnahmen kann dem Kind Angst nehmen.
%   \item \EE{Reden} Sie mit dem Kind und sprechen sie
%   es direkt und mit seinem Namen an! Vermeiden  Sie es über
%   das Kind in der 3. Person zu sprechen, damit das Kind
%   nicht das Gefühl  bekommt, dass hinter seinem Rücken
%   getuschelt wird. Dies würde dem Kind unnötig Angst machen!
  \item \EE{Direktes Ansprechen}: Kinder sind immer direkt anzusprechen!
  Unterhält man sich ständig mit den Eltern über das Kind, so kann es das Gefühl
  haben, dass hinter seinem Rücken getuschelt wird. Dies würde dem Kind unnötig
  Angst machen!
%   \item Geben Sie dem Kind ein Spielzeug in die Hand!
  \item \EE{Spielzeug} kann Kinder (von der Angst) ablenken. 
%   \item \EE{Bezugsperson}: Trennen Sie das Kind niemals von
%   der Bezugsperson! Hat das Kind Angst vor einer bestimmten
%   Untersuchung oder einer Maßnahme, zeigen Sie diese zuerst
%   an der Bezugsperson oder an einem Stofftier vor!
  \item \EE{Bezugsperson}: Kinder brauchen die Nähe ihrer Eltern
  (Bezugspersonen)! Kinder daher niemals von diesen trennen. Außerdem kann man
  bevorstehende Untersuchungen oder Maßnahmen an Bezugspersonen (oder
  Stofftieren) vorzeigen, sodass dem Kind Angst genommen wird.
%  \item Loben Sie das Kind nach der Untersuchung bzw. nach der Maßnahme!
  \item \EE{Lob} hilft ebenfalls weiter!
%   \item \EE{Zeigen lassen}: Da Kinder ihre Beschwerden oft nicht
%   eindeutig beschreiben können, lassen Sie sich z.\,B. den
%   Ort der Schmerzen zeigen. Auf die \E{Körperhaltung} u.\,ä.
%   äußere Anzeichen muss besonders
%   geachtet werden!
  \item \EE{Zeigen lassen}: Da Kinder ihre Beschwerden oft nicht
  eindeutig beschreiben können, ist es sinnvoll sich z.\,B. den
  Ort der Schmerzen zeigen zu lassen. Außerdem muss besonders auf die
  \E{Körperhaltung} u.\,ä.
  äußere Anzeichen geachtet werden!
\end{itemize}
}{
\begin{itemize}
  \item \E{Aufrichtigkeit} und \E{Aufklärung} gelten auch
  für Kinder
  \item Aufrichtigkeit
  \item Aufklärung
  \item Erklären
  \item Reden
  \item Bezugsperson
  \item Sachen \E{zeigen} lassen
\end{itemize}
}{}{}{}



\section{Anamnese}


Als \T{Anamnese} bezeichnet man das Erheben der
Krankengeschichte. Dies schließt sowohl frühere und
chronische Erkrankungen, Operationen und Verletzungen, als
auch auch das Erheben von Details der gegenwärtigen Beschwerden ein. 
Wichtigstes
Instrument zur Erhebung der Anamnese ist das
\E{Patientengespräch}, aber auch eine \I{Fremdanamnese}
(Informationserhebung bei Angehörigen oder anderen
Beteiligten) oder Dokumentationen anderer
Gesundheitseinrichtungen (\E{Patientenbriefe}, \ldots) sind
wichtige Hilfsmittel.



\Layout{27}
{}
{}
{}
{./images/xkcd/desert_island}
{Stille Wasser sind tief \ldots}
{xkcd.com, CC-BY-NC-2.5}


\subsection{Struktur durch \enquote*{SAMPLER}}
\index{SAMPLER}\index{S-KAMEL}

\TextSlide{\TxSAMPLER}{%
Bei einem Notfalleinsatz soll das Anamnesegespräch
gleichzeitig alle wichtigen Informationen zu Tage fördern,
andererseits aber effizient und zeitsparend durchgeführt
werden. Um diese Gegensätze miteinander vereinen zu können 
soll es strukturiert durchgeführt werden. Normalerweise
enthält es die folgenden wichtigen Punkte:

\begin{itemize}
    \item \EE{\color{DefaultA}S}ymptome und Schmerzen
    \item \EE{\color{DefaultA}A}llergien
    \item \EE{\color{DefaultA}M}edikamente
    \item \EE{\color{DefaultA}P}atientengeschichte (Vor- und
    Grunderkrankungen, Operationen, \ldots)
    \item \EE{\color{DefaultA}L}etzte \ldots
    \item \EE{\color{DefaultA}E}reignisse vor
    Notfalleintritt
    \item \EE{\color{DefaultA}R}isikofaktoren
\end{itemize}

In Anlehnung an die Anfangsbuchstaben wird dieses
Abfrageschema \T{SAMPLER}\footnote{Sample \Lang{engl.}:
Probe}\footnote{Manchmal wird dieses Schema 
\I{S-KAMEL} genannt, für Symptome, Krankheiten, Allergien,
Medikamente, Ereignisse und Letztes.} bezeichnet.
Ein derartiges Schema kann im hektischen
Notfall Sicherheit geben. Das Abarbeiten der Fragen
stellt sicher, dass keine wichtigen Informationen
vergessen werden. In der Praxis kann die Reihenfolge der
Fragen je nach Situation abweichen.

}{%
\begin{itemize}
    \item \EE{\color{DefaultA}S}ymptome und Schmerzen
    \item \EE{\color{DefaultA}A}llergien
    \item \EE{\color{DefaultA}M}edikamente
    \item \EE{\color{DefaultA}P}atientengeschichte 
    \item \EE{\color{DefaultA}L}etzte \ldots
    \item \EE{\color{DefaultA}E}reignisse vor
    Notfalleintritt
    \item \EE{\color{DefaultA}R}isikofaktoren
\end{itemize}
}{}{}{}




\TextSlide{Symptome und Schmerzen}{%
\index{Schmerz}Die Einstiegsfrage \Speech{\enquote{Was
kann ich für Sie tun?}} weist oft
bereits auf das Haupt- bzw. Leitsymptom hin. Ist dies
nicht der Fall, so muss als erstes nach dem Hauptsymptom gefragt werden:

\begin{quote}
    \Speech{\enquote{Welche Beschwerden haben Sie?}}
    
    \Speech{\enquote{Was können wir für Sie tun?}}
\end{quote}
Außerdem muss immer auch speziell nach Schmerzen gefragt werden!
\begin{quote}
    \Speech{\enquote{Haben Sie (dabei) Schmerzen?}}
\end{quote}

\subparagraph{Fragen zu Symptomen -- OPQRST}
Gibt der Patient Symptome an, müssen diese weiter abgeklärt
werden, hier am Beispiel Schmerzen:
\index{Schmerz!Arten}

\begin{itemize}
  \item \EE{Lokalisation}: \Speech{\enquote{Wo sind die
  Schmerzen?}}
  \item
  \EE{Qualität}: \Speech{\enquote{Wie ist der Schmerz? Bohrend,
  stechend, brennend, ziehend, krampfartig, auf- und
  abschwellend?}}
  \item \EE{Ausstrahlung}: \Speech{\enquote{Strahlt der Schmerz
  irgendwohin aus?}}
  \item \EE{Zeitdauer / Beginn}: \Speech{\enquote{Seit
  wann haben Sie die Schmerzen? Wie lange haben Sie die Schmerzen schon?
  Haben Sie so einen Schmerz schon einmal gehabt? Haben die
  Schmerzen plötzlich oder langsam eingesetzt?}}
  \item \EE{Geschehnisse zum Zeitpunkt des Auftretens}:
  \Speech{\enquote{Was haben Sie gemacht als es begonnen
  hat?}}
  \item \EE{Provokation / Palliation}:
  \Speech{\enquote{Was macht den Schmerz besser oder schlechter?}}
  \item \EE{Stärke}: \Speech{\enquote{Wie stark ist der
  Schmerz?}}
  Setzen Sie die Aussage des Patienten immer in Relation
  zu seinem Verhalten, seinem Gesichtsausdruck und zu
  seiner Körperhaltung! Einem Patienten mit starken
  Schmerzen sieht man dies oft, aber nicht immer
  an!\ZFootnote{\E{Stärke des Schmerzes}: Einige Autoren
  empfehlen, den Schmerz anhand einer Skala von
  \VonBis{0}{10} beurteilen zu lassen.
  % %   Beurteilung: ex STEUER Review 2011-07-11
  %     Eine solche Beurteilung ist jedoch nur
  %     im Verlauf sinnvoll, da der Patient im akuten
  %     Stadium keinen Anhaltspunkt hat, um die Stärke der
  %     Schmerzen zuverlässig in Relation zu setzen und
  %     einzustufen.
  }
\end{itemize}
\Z{Diese
Fragenabfolge wird in Anlehnung auf die Benennung der
EKG-Zacken als \T{OPQRST} bezeichnet\ZFootnote{\I{OPQRST}: Onset
(Geschehnisse zum Zeitpunkt des Auftretens), Provocation or
Palliation (Was macht es besser bzw. schlechter), Quality (Schmerzqualität,
-art), Region and Radiation (Lokalisation und
Ausstrahlung), Severity (Stärke), Time (zeitlicher
Verlauf)}}.

}{
% \begin{itemize}
%   \item Wo?
%   \item Ausstrahlung?
%   \item Seit wann? Wie lange schon?
%   \item Plötzlicher oder langsamer Beginn?
%   \item Verlauf?
%   \item Wie?
%   \item Wie stark?
%   \item Was provoziert den Schmerz?
% \end{itemize}
% 
\TabHeading{OPQRST}

\begin{itemize}
  \item \EE{Lokalisation}: \Speech{\enquote{Wo sind die
  Schmerzen?}}
  \item
  \EE{Qualität}: \Speech{\enquote{Wie ist der Schmerz? Bohrend,
  stechend, brennend, ziehend, krampfartig, auf- und
  abschwellend?}}
  \item \EE{Ausstrahlung}: \Speech{\enquote{Strahlt der Schmerz
  irgendwohin aus?}}
  \item \EE{Zeitdauer / Beginn}: \Speech{\enquote{Seit
  wann haben Sie die Schmerzen? Wie lange haben Sie die Schmerzen schon?
  Haben Sie so einen Schmerz schon einmal gehabt? Haben die
  Schmerzen plötzlich oder langsam eingesetzt?}}
  \item \EE{Geschehnisse zum Zeitpunkt des Auftretens}:
  \Speech{\enquote{Was haben Sie gemacht als es begonnen
  hat?}}
  \item \EE{Provokation / Palliation}:
  \Speech{\enquote{Was macht den Schmerz besser oder schlechter?}}
  \item \EE{Stärke}: \Speech{\enquote{Wie stark ist der
  Schmerz?}}
\end{itemize}
}{}{}{}





\TextSlide
{Besonderheit: Kolikartiger Schmerz}
{\label{topic:kolik}
\index{Kolik!Schmerzqualität}
\index{Schmerz!kolikartiger}

Ein \T{Kolikartiger Schmerz} ist \EE{an-}
und \EE{abschwellend}.
Er wird oft als krampfhaft und sehr stark beschrieben. Er
kommt häufig bei Verstopfung der Gallen- und Harnwege durch
Steine vor (\FullPrettyRef{topic:gallenkolik},
\FullPrettyRef{topic:nierenkolik}).}
{
\begin{compactitem}
  \item Auf- und abschwellend
  \item Krampfhaft
\end{compactitem}
}
{}
{}
{}



\Layout{2}{Allergien}{%
Auf die Frage nach möglichen Allergien wird häufig
vergessen, ein Informationsmangel kann aber fatale
Auswirkungen haben. Es muss daher immer nach Allergien,
in einer zweiten Frage sogar speziell nach
Medikamentenunverträglichkeiten gefragt werden! Bedenke: Ein
Notfallpatient kann bewusstlos werden und vielleicht zu
einem späteren Zeitpunkt nicht mehr befragt werden!
}{}{}{}




\TextSlide{Medikamente}{%
Die vom Patienten eingenommenen Medikamente lassen
Rückschlüsse auf seine Grunderkrankungen zu und sind im Falle
eines längeren Spitalsaufenthalts eine wertvolle Information.
Auch die Regelmäßigkeit oder Unregelmäßigkeit der Einnahme
kann Ihnen Anhaltspunkte für das Finden Ihrer
Verdachtsdiagnose bzw. der Umstände des Patienten liefern.
Angaben zu Medikamenten sind für die nachbehandelnde
Einrichtung zu dokumentieren. Es muss immer nach
Medikamentenlisten o.\,ä.  gefragt werden, diese sind für
die weiterbehandelnde Stelle mitzunehmen! Sind keine Listen
vorhanden, so  müssen die entsprechenden
Medikamentenpackungen des Patienten gesichert werden.

\Fazit{Frage immer nach Medikamentenlisten!}

\subparagraph{Fragen zu Medikamenten:}
\begin{quote}
    \emph{\enquote{Welche Medikamente nehmen Sie? Haben
    Sie eine Liste?}}
    
    \emph{\enquote{Wogegen nehmen Sie diese Medikamente?}}
    
    \emph{\enquote{Nehmen Sie Ihre Medikamente regelmäßig?}}
    
    \emph{\enquote{Wann haben Sie Ihre Medikamente zuletzt genommen?}}
    
\end{quote}
}{
\begin{itemize}
  \item Welche?
  \item Wogegen?
  \item Regelmäßige Einnahme?
  \item Wann zuletzt genommen?
  \item \EE{Medikamentenliste}
\end{itemize}
}{}{}{}


\Layout{57}{}
{}
{}
{\scriptsize\begin{tabulary}{\linewidth}{LLLLLL}
\toprule\addlinespace
\TabHeading{(Marken-) Name}
				& \TabHeading{(Grund-) Erkrankungen}
									& \TabHeading{Anmerkungen}	
													& \TabHeading{Andere Namen / Generika}		
																					& \TabHeading{Wirkstoff} 	
																									& \TabHeading{Wirkung} 			
\\\addlinespace\midrule\addlinespace
\noindent\Ca{Marcoumar{\texttrademark}}
				& z.\,B. Herzrhytmusstörungen
									& Massiv gesteigerte \EE{Blutungsneigung}
							 						& 								&				& Blutgerinnung ↓ 						
\\\addlinespace
Aspirin{\texttrademark} 		
				& Schmerzen, gripp. Infekte
									& Gesteigerte Blutungsneigung, schädigt Magenschleimhaut
													& \Z{unzählige, z.\,B.: ASS-Genericon{\texttrademark}, ASS-Ratiopharm{\texttrademark}, \ldots }
																					& \Z{Acetyl-Salicylsäure (ASS)} 
																									& Gerinnungshemmend\Z{, schmerzhemmend, entzündungshemmend, fiebersenkend} 
\\\addlinespace
Mexalen{\texttrademark} 		
				& Schmerzen, gripp. Infekte
									& Lebertoxisch, Vergiftungen häufig
													&								& Paracetamol 	& \Z{Fiebersenkend, entzündungshemmend, schmerzhemmend}
																						
\\\addlinespace
Digitalis{\texttrademark} 		
				& Herzinsuffizienz, Herzrhythmusstörungen
									& Über-/Unterdosierungen häufig, kann Herzrhythmusstörungen verursachen
													& 								&				& \Z{Herzkraftsteigernd, Herzrhythmus-stabilisierend}
\\\addlinespace
Diverse; z.\,B.: Voltaren{\texttrademark}, Aspirin{\texttrademark}, \ldots 
				& Schmerzen, gripp. Infekte					
									& Schädigung der Magenschleimhaut
													& Diverse 						& \Z{Nicht-steroidale Antirheumatika (NSAR), z.\,B.: Diclofenac, ASS, \ldots}
																									& \Z{Fiebersenkend, entzündungshemmend, schmerzhemmend} 												
\\\addlinespace
Insulin		& Diabetes mellitus	& Hypoglykämie; Hyperglykämie bei Nichteinnahme				
													& Diverse, verschiedene lang- und kurzwirksame Präparate verfügbar							
																					& Insulin		& Blutzucker ↓												
\\\addlinespace
ThromboASS{\texttrademark}		
			& Z.\,n	Herzinfarkt, Z.\,n. Insult
								&  Anamnese!
													& 							
																					& \Z{Acetyl-Salicylsäure (ASS)}		
																									& \Z{Gerinnungshemmend}											
\\\addlinespace
Plavix{\texttrademark}		
			& Z.\,n	Herzinfarkt, Z.\,n. Insult
								&  Anamnese!
													& 							
																					& \Z{Clopidogrel}		
																									& \Z{Gerinnungshemmend}
\\\addlinespace
Nitroglycerin		
			& Koronare Herzkrankheit
								&  Erhältlich u.\,a. als Spray, Kapseln und Pflaster
													& 							
																					& Nitrolglycerin	
																									& Erweiterung der Herzkranzgefäße\par Achtung:
																									Blutdruck ↓
																					
\\\addlinespace\bottomrule
\end{tabulary}

{\noindent\footnotesize Andere ähnliche Präparate, welche die Gerinnung wie Marcoumar beeinflussen, wären z.\,B. Pradaxa{\texttrademark} oder Sintrom{\texttrademark}}}
{Besonders häufige bzw. wichtige Medikamente}
{}

Notizen: \TakeNotesLarge



\TextSlide
{Patientengeschichte}
{Notfälle ereignen sich manchmal aufgrund der akuten
Verschlechterung einer bestehenden Vorerkrankung. Außerdem
hilft die Beurteilung der Messwerte, wenn man über die
Grunderkrankungen des Patienten Bescheid weiß. Zum Beispiel
kann die Ursache eines hohen Blutdrucks eine bereits
bestehende Erkrankung, aber auch ein Notfallgeschehen sein.
Neben der Befragung können auch Dokumentationen von anderem
Gesundheitsfachpersonal (Spitals- bzw. Arztbriefe, Befunde,
Pflegedokumentation, etc.) Auskunft über bestehende
Vorerkrankungen liefern.

% OP, Dauer-erkr. \ldots

Allgemeine Fragen,
wie z.\,B.
\begin{quote}
    \Speech{\enquote{Leiden Sie unter irgendwelchen
    Krankheiten?}}
    
    \Speech{\enquote{Bestehen bei Ihnen Vorerkrankungen?}}
\end{quote}
können anschließend durch Symptom-spezifische Fragen ergänzt
werden:
\begin{quote}
    \Speech{\enquote{Hatten Sie schon einmal einen
    Herzinfarkt?}}
    
    \Speech{\enquote{Leiden Sie unter hohem Blutdruck?}}
    
    \Speech{\enquote{Sind Sie zuckerkrank?}}
\end{quote}}
{
\begin{itemize}
  \item Grund- und Vorerkrankungen?
  \item Behinderungen?
  \item Operationen?
  \item \ldots
\end{itemize}
}
{}
{}
{}


\TextSlide{Letzte \ldots}{%
Dieser Punkt des Abfrageschemas bezieht sich auf letzte
Ereignisse, welche mit dem Notfall in Zusammenhang stehen
könnten. Damit Sie hier die richtigen Fragen stellen können,
ist ein fundiertes Fachwissen notwendig. Abhängig von den
Beschwerden können verschiedene Fragen zielführend sein:
\begin{quote}
    \emph{\enquote{Wann haben Sie zuletzt etwas
    gegessen/getrunken?}} (Ist der Patient nüchtern?)
    
    \emph{\enquote{Wann hatten Sie zuletzt Stuhlgang?}}
    (weiterführende Frage: Wie hat der Stuhl ausgesehen?)
    
    \emph{\enquote{Wann war die letzte Regelblutung?}}
    (Besteht die Möglichkeit einer Schwangerschaft?)
    
    \emph{\enquote{Wann waren Sie zuletzt beim Arzt bzw. im
    Spital?}}
\end{quote}
}{
\begin{itemize}
  \item Letzte Mahlzeit?
  \item Letzter Stuhlgang?
  \item Letzte Regelblutung?
  \item Letzter Spitalsaufenthalt?
  \item Letzter Arztbesuch?
  \item Letzte Blutzuckerkontrolle?
  \item \ldots
\end{itemize}
}{}{}{}




\Layout{2}{Ereignisse vor Notfalleintritt}{%
Oft ereignen sich vor einem Notfall Schlüsselereignisse,
welche den Notfall herbeiführen oder diesen auslösen.
Fragen Sie daher immer nach möglichen Besonderheiten oder
Tätigkeiten des Patienten vor dem Notfall:
\begin{quote}
    \Speech{\enquote{Was haben Sie vor Beginn der Symptome
    gemacht?}}
    
    \Speech{\enquote{Haben Sie sich vor Beginn der Symptome
    besonders angestrengt?}}
    
    \ldots{}
\end{quote}
In dieser Kategorie kann es auch hilfreich sein, wenn Sie
Angehörige oder Zeugen befragen:
\begin{itemize}
    \item Hat der Patient vor dem Notfall etwas eingenommen
    (Überdosierung)?
    \item Hat der Patient gekrampft bevor er bewusstlos liegen geblieben ist?
    \item Ist der Patient vor dem Eintreten der Symptome gestürzt?
    \item u.\,v.\,m.
\end{itemize}
}{

}{}{}{}

\TextSlide{Risikofaktoren}{%
Viele Erkrankungen werden durch das Vorhandensein von
Risikofaktoren begünstigt.
Das Vorliegen von Riskofaktoren bedeutet natürlich nicht,
dass eine bestimme Erkrankung vorliegt, die Kenntnis über
Risikofaktoren kann aber in der Diagnosefindung hilfreich
sein und das Gesamtbild des Patienten vervollständigen.
Beispiele für Riskofaktoren wären:
\begin{itemize}
  \item Alter
  \item Geschlecht
  \item Immobilisierung, Bettlägrigkeit (Thrombosen,
  Lungenembolie, Pneumonie, \ldots)
  \item Rauchen (Lungenkrebs, COPD, Lungenembolie, \ldots)
  \item Diabetes mellitus (Herz-Kreislauf-Erkrankungen:
  Koronare Herzkrankheit, Herzinfarkt, Schlaganfall,
  periphere arterielle Verschlusskrankheit (pAVK),
  Wundheilungsstörungen, \ldots; Niereninsuffizienz, \ldots)
  \item Arterielle Hypertonie (Herz-Kreislauf-Erkrankungen, s.\,o.)
  \item Hormonelle Verhütung (\enquote*{Pille}; Lungenembolie, \ldots)
\end{itemize}
Die Frage nach Risikofaktoren ist abhängig vom Einzelfall,
oft sind wichtige Risikofaktoren auch ohne Nachfragen
erkennbar (Bettlägrigkeit, Alter, Geschlecht \ldots). Es ist
wichtig, sich die gefundenen Riskofaktoren bei der Bewertung
des Patienten zu vergegenwärtigen.
}{
\begin{itemize}
  \item Begünstigen Erkrankungen
  \item Z.\,B.:
  \begin{itemize}
    \item Alter
    \item Geschlecht
    \item Immobilisierung, Bettlägrigkeit
    \item Rauchen
    \item Diabetes mellitus
    \item Arterielle Hypertonie
    \item Hormonelle Verhütung
  \end{itemize}
\end{itemize}
}{}{}{}


\clearpage
\section{Dokumentation von anderen Gesundheitseinrichtungen und -berufen}
\label{topic:dokumentation-andere-gesundheitsberufe}

\subsection{Arzt- und Entlassungsbriefe}

Nach einer stationären Aufnahme in einer
Gesundheitseinrichtung wird ein Arzt- bzw. Entlassungsbrief
erstellt. In
diesem Dokument finden sich wichtige Informationen bezüglich
der früheren Krankheitsgeschichte des Patienten, der
Medikation sowie dem Behandlungsverlauf in der jeweiligen
Einrichtung. Eventuell ist auch ein Pflegebericht beigelegt.
Der Brief ist somit eine ideale \E{Ergänzung} des
Anamnesegespräches (aber kann es \E{nicht} ersetzen!).

In jedem Fall ist zu überprüfen, ob es sich tatsächlich um
den Brief des Patienten handelt, und ob die Angaben darin
noch aktuell sind (insbesonders Angaben zu Diagnosen und
Medikation können sehr rasch veralten).

% \bigskip
% \clearpage
\subsubsection{Beispiel: Patient mit schlechtem
Allgemeinzustand}

\hrule

\vspace{1em}

\begin{center}\sffamily
 \textbf{Krankenanstalt der Barmherzigkeit} 
 \\II. Med. Abteilung 
 
Vorstand: Prim. Univ.-Doz. Dr. Vicco von Bülow
 
{\scriptsize Hartlgasse 16a, A-1240 Wien
\\
Tel.: +43\,555\,328\,745\,-\,0 ~~~ Fax: +43\,555\,328\,745\,-\,100}
\end{center}

{\fontfamily{lmtt}\fontseries{lc}\selectfont                                                                                               

\vspace{1em}
                                                                                                                            


 Peter Zapfl\\
Diesseitsweg 13a/12/5/34 \\
A-1147 Wien
\vspace{1.5em}

\begin{flushright}
Wien, am 25.07.2009
\end{flushright}

  Wir berichten über den Aufenthalt unseres Patienten Herrn Peter ZAPFL, geb.
  am 13.12.1936, welcher vom 30.03.2009 bis 17.04.2009 an unserer Abteilung
  stationär aufgenommen war.

{\bfseries   Aus der Anamnese:}

\begin{quote}
  Kindheit: \MarginNo{\Z{AE: Appendeketomie, Entfernung des
  Blinddarms; TE: Tonsillektomie, Entfernung der
  Mandeln.}}AE, sowie TE, WK II: längerer Lazarettaufenthalt
  wegen Typhus.
  
  Seit 2003 sind ein bis dato nicht insulinpflichtiger
  Diabetes mellitus, sowie eine art. Hypertonie bekannt.
  
  2005 Aufenthalt an unserer Abteilung wg. eines im
  \MarginNo{\Z{CT:
  Computertomographie.}}CT nachgewiesenen rechtshirnigen
  ischämischen Insultes, anschließend Rehabaufenthalt in Bad
  Schallerbach.
  
  Zuletzt lag der Patient vom 21.11.2008 bis 23.12.2008
  wegen eines nicht insulinpflichtigen, entgleisten Diabetes
  mellitus, bei Z.~n. ischämischen zerebralen Insult 2005
  mit Halbseitenzeichen links und arterieller Hypertonie an
  unserer Abteilung.
  
  Zwischenzeitlich keine Auffälligkeiten.
  
  Alkohol: verneint. Nikotin: verneint. Koffein: fallweise.
  Harn, Stuhl: unauff.
  Schlaf: unauff.
\end{quote}


{\bfseries Bisherige Med.:} 

\begin{quote}
  \begin{tabular}{lll}
Renitec 10 mg &&1-1-1
\\
Diabetex 850 mg &&1-1-1
\\
Diab Diät 12 BE &&2-1-4-2-2-1
\\
Thromboass 100 mg &&1-0-0
\\
\end{tabular}
\end{quote}





 Die nunmehrige Aufnahme erfolgte mit dem ASB, der wegen zunehmender
 Verschlechterung des AZ des Patienten vom Wohnungsnachbarn
 berufen wurde.

{\bfseries Auffälliges im Aufnahmestatus:}

\begin{quote}
 Deutlich herabgesetzter \MarginNo{AZ: Allgemeinzustand\Z{,
 EZ: Ernährungszustand}.}AZ, noch ausreichender EZ, zeitlich
 und örtlich teilorientiert, motorische Unruhe. Aktive
 Beweglichkeit der re. \MarginNo{OE: Obere Extremität, UE:
 Untere Extremität.}OE und UE, beinbetonte spastische
 Halbseitenlähmung links mit gesteigerten Sehnenreflexen
 links bei Z.\,n. re-Hirn-Insult, mäßige
 Artikulationsstörungen.
 
 Haut trocken, \MarginNo{\Z{Turgor: vom Flüssigkeitsgehalt
 abhängiger Spannungszustand der Haut.
 \nocite{Pschy259}}}Turgor herabgesetzt, kein Dekubitus
 nachweisbar.
 \end{quote}

{\bfseries Befunde:}
\begin{quote}
Labor siehe Beilage.

\Margin{\Z{C/P: \enquote{Cor-Pulmo}, Thoraxröntgen.}}C/P.: Keine
wesentliche Änderung gegenüber Vorbefund von 12/2008.

Schädel CT: Keine wesentliche Änderung gegenüber Letztbefund
von 12/2008, Feinbefund siehe siehe Beilage.

Augen: Erstgradige hypertone Veränderungen der
\MarginNo{\Z{Retina: Netzhaut.}}Retinagefässe, kein Hinweis
für \MarginNo{\Z{Diabetische Retinopathie: Schädigung der
Netzhautgefäße im Rahmen des Diabetes mellitus.}}diabetische
Retinopathie.

Neuro: Anamnestisch Z.\,n. Insult 2005 im Strombereich der
A. cerebri ant., mit beinbetonter spastischer Halbseitenlähmung
li und gesteigerten \MarginNo{\Z{BSR, TSR: Sehnenreflexe des
Bizeps und Trizeps.}}BSR und TSR li sowie stark gesteigerten
\MarginNo{\Z{PSR: Patellarsehnenreflex. Babinski-Reflex:
krankhafter Reflex bei Schädigung des Rückenmarks oder
Gehirns.}}PSR li., Babinski li. pos., Trömner u. Knips
beidseits neg., kein Hinweis auf Re-Insult, weiterhin
physi\-ka\-lisch- thera\-peutische und logopädische
Maßnahmen empfohlen. Der durchgeführte Schlucktest ergab
keinen Hinweis auf eine relevante Schluckstörung.

Uro: Prost. rect. dig. unauffällig.

EKG: \MarginNo{\Z{Sinusrhythmus, Herzfrequenz 74,
Linkstyp.}}SR 74, LT, altersentsprechender Kurvenverlauf.

Der Blutdruck war bei mehrfachen Kontrollen gegen Ende des
Aufenthaltes normoton, zuvor im hypotonen Bereich.
\end{quote}




{\bfseries Epikrise:}

\begin{quote}
  Nach Behebung des bei Aufnahme bestehenden
  Flüssigkeitsdefizites mit geeigneter Infusionstherapie,
  besserte sich der AZ des Patienten zusehends. Mit in der
  Folge durchgeführten physi\-ka\-lisch- therap. Maßnahmen
  konnte der Mobilisisationsgrad des Pat. günstig
  beeinflusst werden.
  
  Der Pat. wird bis zum selbständigen Gehen mit Rollator
  mobilisiert entlassen. Die zu Aufenthaltsbeginn durchwegs
  im hypotonen Bereich gemessenen RR-Werte haben sich in
  mehrfachen Kontrollen unter dem angegeben antihypertonen
  Regime normalisiert. Ebenso hat sich die ursprünglich
  unbefriedigend vorgefundene diabetische Stoffwechsellage,
  bei konsequenter Einhaltung des vorgeschlagenen Regimes im
  therapeutischen Bereich stabilisiert.
  
  Im Einverständnis mit dem Pat. und dessen Tochter wurde
  ein Platz im Bernhard-Prischl-Heim zur Gewährleistung
  einer regelmäßigen und auch ausreichenden
  Flüssigkeitszufuhr, sowie einer gesicherten Diabetes-
  Diät, organisiert. Eine Kontrolle der vor etwa einem Jahr
  durchgeführten \MarginNo{\Z{Carotissonographie:
  Ultraschall der Halsschlagader.}}Carotissonographie ist
  angezeigt. Ein entspr. Termin wurde am 27.12.2009  12,30
  Uhr an der Angiolog Ambulanz d. KA Mitschka-Stiftung
  fixiert. Eine Befund\-be\-sprech\-ung ist nach
  Terminvereinbarung an h.o. Int. Ambulanz möglich.
\end{quote}

\MarginNo{\Z{}}
{\bfseries Hauptdiagnose:}

\begin{quote}
\begin{tabular}{ll}
Hypohydratation & E86.0
\\
\end{tabular}
\end{quote}

{\bfseries Nebendiagnosen:}

\begin{quote}
\begin{tabular}{ll}
entgleister. Diabetes mellitus. (NIDDM) & E12.31, E12.40,
E12.72
\\
Art. Hypertonie & I10.1
\\
Z.n. Zerebralem Insult mit HSS li 2005& I63.3
\\
Spastische Hemiparese li, beinbetont&I69.3, G81.1
\\
\end{tabular}
\end{quote}



{\bfseries Therapie:}

\begin{quote}
\begin{tabular}{lll}
Renistad Tab. 5,0 mg && 1-1-1 
\\
Glucophage 850 mg Tab && 1-1-1 
\\
Diab. Diät 12 BE&&(2-1-4-2-2-1) 
\\
Thromboass Tab.&&1-0-0 
\\
tgl. Flktszufuhr nicht unter 2 l&&
\end{tabular}
\end{quote}




{\bfseries Kontrollen:}

\begin{quote}
  NFP, Harnbefund, Elektrolyte, RR, BZ, (Patient. führt seit
  Jahren Selbstkontrollen durch und führt darüber
  Aufzeichnungen)) weiterhin neurolog. und ophthalmolog.
  Kontrollen.
\end{quote}









Mit kollegialen Grüßen\vspace{2em}\\ Dr. Adolf Schrammel, Stationsarzt

\Hand{Schrammel}

vidit OA Univ.Doz. Dr. Alois Schremser
}

\vspace{2em}

\hrule




% \clearpage
% \subsubsection{Beispiel: Patient nach Operation}
% 
% \hrule
% 
% \vspace{1em}
% 
% \begin{center}\sffamily
%  \textbf{Landesklinikum Hinter a. d. Tupfing}
%  \\Abteilung für Chirurgie
% 
% Vorstand: Prim. Dr. Viktor Hauer
% 
% {\scriptsize Spitalgasse 2, A-3212 Hinter/Tupfing
% \\
% Tel.: +43\,555\,328\,745\,-\,0 ~~~ Fax: +43\,555\,328\,745\,-\,100}
% \end{center}
% 
% {%\fontfamily{lmtt}\fontseries{lc}\selectfont
% \sffamily
% 
% \vspace{1em}
% 
% 
% 
%  Walter Mayer\\
% Tako-Tsubo-Weg 14a/11/25 \\
% A-1147 Wien
% \vspace{1.5em}
% 
% \begin{flushright}
% Hinter/Tupfing, am 25.08.2012
% \end{flushright}
% 
%   Wir berichten über den Aufenthalt unseres Patienten Herrn Walter M\,A\,Y\,E\,R, geb.
%   am 01.01.1933, welcher vom 23.08.2012 bis 25.08.2012 an
%   unserer Abteilung stationär aufgenommen war.
% 
% 
% \begin{description}
%   \item[Grund der Aufnahme:] Obstipation, V. a. Ileus
%   \item[Wichtiges zur Aufnahme und Anamnese:] Herr Mayer war
%   bereits wiederholt h.\,o. wegen Obstipation bei bekanntem
%   N. recti in Behandlung. Eine Rektumamputation mit Anus
%   praeter wurde vom Patienten bisher abgelehnt. Zuletzt erfolgte eine palliative
%   Tumorabtragung. Numehr
%   erfolgt die stationäre Aufnahme über die zentrale
%   Notfallambulanz wegen Stuhlverhalt, angeblich seit 2 Wochen.
% \end{description}
% 
% 
% \begin{description}
%   \item[Status präsens:] Reduzierter AZ, reduzierter EZ, \\
%   Cor: rhythmisch, tachykard, rein; Pulmo: vesikuläres
%   Atemgeräusch beidseits, sonorer Klopfschall, Sprechdyspnoe
%   Abdomen: weich, spärliche Darmgeräusche, kein
%   Druckschmerz, Stuhlwalzen\\
%   Neurologisch; grob neurologisch unauffällig
%   \item[Auffälliges im Aufnahmelabor:] siehe Beilage
%   \item[EKG:] VH-Flimmern, 80/min, LAHB
%   \item[Cor-Pulmo-Röntgen:] Keine Infiltrate, keine
%   Stauungszeichen
%   \item[Digital-rektale Untersuchung:] Resistenz ca. 4 cm ab
%   ano, Ampulle stuhlfrei
% \end{description}
% 
% \paragraph{Verlauf:}
% Der Patient wird umgehend auf der Herzüberwachungsstation
% aufgenommen und überwacht.
% 
% \paragraph{Diagnosen}~
% \par
% 
% \begin{tabulary}{\linewidth}{rL}
% --.-
% &
% \textbf{Tachykardes Vorhofflattern}
% \\\addlinespace
% & Z.n. Vorhofflimmern und medikamentöser Cardioversion
% 12/2011
% \\
% &Z.n. Pulmonalvenenisolation 3/2012
% \\
% &Koronarsklerose
% \\
% & Ausschluss sign. Stenose 3/2012
% \\
% I10.1
% &
% Art. Hypertonie
% \\
% &Z.n. AE, TE
% \end{tabulary}
% 
% \paragraph{Therapie bei Entlassung}~
% 
% \begin{tabulary}{\linewidth}{Ll}
% Sedacoron 200mg & 1-0-0
% \\
% Concor 5mg & 1-0-0
% \\
% ThromboASS 100mg & 1-0-0
% \\
% Simvastatin 20mg & 0-0-1
% \\
% Marcoumar & lt. Pass
% \end{tabulary}
% 
% 
% \paragraph{Weiteres Procedere:} ~
% 
% 
% 
% 
% \bigbreak
% 
% 
% 
% \noindent\begin{tabularx}{\linewidth}{XXX}
% Oberarzt & & Sekundararzt
% \\
% \textbf{OA Dr. R. Glocke}
% &
% &
% Sek.-Dr. H. Harrer
% \\\addlinespace
% \\
% \Hand{Glocke}
% 
% \end{tabularx}
% \vspace{2em}
% 
% \hrule
% }





% 
% 
% 
% 
% 
% 
% 
% 
% 
% 


% \clearpage
% \subsubsection{Beispiel: Patient mit tachykarder
% Rhythmusstörung}
% 
% \hrule
% 
% \vspace{1em}
% 
% \begin{center}\sffamily
%  \textbf{Landesklinikum Hinter a. d. Tupfing}
%  \\Abteilung für interne Medizin
% 
% Vorstand: Prim. Dr. Viktor Hauer
% 
% {\scriptsize Spitalgasse 2, A-3212 Hinter/Tupfing
% \\
% Tel.: +43\,555\,328\,745\,-\,0 ~~~ Fax: +43\,555\,328\,745\,-\,100}
% \end{center}
% 
% {%\fontfamily{lmtt}\fontseries{lc}\selectfont
% \sffamily
% 
% \vspace{1em}
% 
% 
% 
%  Walter Mayer\\
% Tako-Tsubo-Weg 14a/11/25 \\
% A-1147 Wien
% \vspace{1.5em}
% 
% \begin{flushright}
% Hinter/Tupfing, am 25.08.2012
% \end{flushright}
% 
%   Wir berichten über den Aufenthalt unseres Patienten Herrn Walter M\,A\,Y\,E\,R, geb.
%   am 01.01.1970, welcher vom 23.08.2012 bis 25.08.2012 an
%   unserer Abteilung stationär aufgenommen war.
% 
% 
% \begin{description}
%   \item[Grund der Aufnahme:] Tachykarde Rhythmusstörung,
%   Schmalkomplextachykardie
%   \item[Wichtiges zur Aufnahme und Anamnese:] Herr Mayer war
%   wiederholt h.\,o. nach Episoden von Herzklopfen und Z.\,n. Kollaps vorstellig. Im Dezember
%   2011 konnte erstmalig ein Vorhofflimmern
%   im EKG dargestellt werden. Seitdem besteht eine orale
%   Antokoagulation mit Marcoumar und eine Rezidivprophylaxe
%   mit Sedacoron und Concor. Eine im LK Krems durchgeführt
%   Koronarangiographie ergab keine mäiggradige
%   Koronarsklerose ohne Hinweis auf eine signifikante
%   Stenosierung. Es wurde ferner im März 2012 eine
%   Pulmonalvenenisolation durchgeführt.
% \end{description}
% 
%   \paragraph{Grund der Aufnahme:} Tachykarde
%   Rhythmusstörung, Schmalkomplextachykardie
%   \paragraph{Wichtiges zur Aufnahme und Anamnese:} Herr
%   Mayer war wiederholt h.\,o. nach Episoden von Herzklopfen und Z.\,n. Kollaps vorstellig. Im Dezember
%   2011 konnte erstmalig ein Vorhofflimmern
%   im EKG dargestellt werden. Seitdem besteht eine orale
%   Antokoagulation mit Marcoumar und eine Rezidivprophylaxe
%   mit Sedacoron und Concor. Eine im LK Krems durchgeführt
%   Koronarangiographie ergab keine mäiggradige
%   Koronarsklerose ohne Hinweis auf eine signifikante
%   Stenosierung. Es wurde ferner im März 2012 eine
%   Pulmonalvenenisolation durchgeführt.
% 
%   Der Patient kommt nun bei Z.\,n. Kollaps und
%   Palpitationen. Im Aufnahme-EKG zeigt sich ein tachykardes
%   Vorhofflattern mit 2:1-Überleitung.
% 
% \begin{description}
%   \item[Status präsens:] Reduzierter AZ, adipöser EZ, \\
%   Cor: rhythmisch, tachykard, rein; Pulmo: vesikuläres
%   Atemgeräusch beidseits, sonorer Klopfschall, Sprechdyspnoe
%   Abdomen: weich, unauff. Darmgeräusche, kein Druckschmerz\\
%   Neurologisch; grob neurologisch unauffällig
%   \item[Auffälliges im Aufnahmelabor:] siehe Beilage
%   \item[EKG:] VH-Flattern, 140/min, LAHB
%   \item[Cor-Pulmo-Röntgen:] Keine Infiltrate, keine
%   Stauungszeichen
%   \item[]
% \end{description}
% 
% \paragraph{Verlauf:}
% Der Patient wird umgehend auf der Herzüberwachungsstation
% aufgenommen und überwacht.
% 
% \paragraph{Diagnosen}~
% \par
% 
% \begin{tabulary}{\linewidth}{rL}
% --.-
% &
% \textbf{Tachykardes Vorhofflattern}
% \\\addlinespace
% & Z.n. Vorhofflimmern und medikamentöser Cardioversion
% 12/2011
% \\
% &Z.n. Pulmonalvenenisolation 3/2012
% \\
% &Koronarsklerose
% \\
% & Ausschluss sign. Stenose 3/2012
% \\
% I10.1
% &
% Art. Hypertonie
% \\
% &Z.n. AE, TE
% \end{tabulary}
% 
% \paragraph{Therapie bei Entlassung}~
% 
% \begin{tabulary}{\linewidth}{Ll}
% Sedacoron 200mg & 1-0-0
% \\
% Concor 5mg & 1-0-0
% \\
% ThromboASS 100mg & 1-0-0
% \\
% Simvastatin 20mg & 0-0-1
% \\
% Marcoumar & lt. Pass
% \end{tabulary}
% 
% 
% \paragraph{Weiteres Procedere:} ~
% 
% 
% 
% 
% \bigbreak
% 
% 
% 
% \noindent\begin{tabularx}{\linewidth}{XXX}
% Oberarzt & & Sekundararzt
% \\
% \textbf{OA Dr. R. Glocke}
% &
% &
% Sek.-Dr. H. Harrer
% \\\addlinespace
% \\
% \Hand{Glocke}
% 
% \end{tabularx}
% \vspace{2em}
% 
% \hrule
% }

\begin{Literary}
\fullcite{DahmerGespraechsfuehrung5}

\fullcite{FamPropAeGeFue2010}

\fullcite{ThunRedenI46}

\fullcite{ArgelanderErstinterview8}
\end{Literary}





\ChapterEnd






